\scanpage{144}
\mspage{89}
% archive sheet number: 143

Donc on a sans larmes les deux foncteurs $\delta,\unsure{\gamma|\sigma}$ comme quasi-inverses l'une de l'autre. Et le lemme page 11 montre même que pour \uline{\unsure{tte|toute}} sous-catégorie pleine $\unsure{\mathcal F|F}$ de $\deleted{\unsure{\mathrm{C}}}\mathcal M_{/S}$, telle que

\[
\underline{\mathrm{Fib}}\,S\subset\mathcal F\subset\mathcal M_{/S},
\]

les inclusions induisent des équivalences sur les localisées $W_S$

\diagram{\begin{tikzcd}[ampersand replacement=\&,column sep=1.25em,row sep=0.4em]
(\underline{\mathrm{Fib}}\,S)W_S^{-1}\arrow[r,-{Bar},"\approx"]\&\mathcal F W_S^{-1}\arrow[r,-{Bar},"\approx"]\&\underline{(\mathcal M_{/S})W_S^{-1}}\arrow[d,equal,"\mathrm{def}"]\\
\&\&\mathrm{Ho}^{\unsure{\mathrm{l.c}|\mathrm{l.e}}}_W(S).
\end{tikzcd}}

Je vais traiter maintenant le cas du type Quillen, pour me convaincre que l'existence du foncteur $C$, et \unsure{les choix particuliers|le choix particulier} qu'on peut en faire, est \unsure{irrelevant|irrelev.nt}. (Ce dernier point est d'ailleurs évident à posteriori, quand on sait que $\delta$ est une équivalence, car le qu.inv. $\gamma$ sera alors unique à isom unique près, et le choix d'un foncteur $C$, et d'un \unsure{hom|lien} fonctoriel $i:\mathrm{id}\unsure{-|\to}C$, n'est qu'une façon parmi d'autres d'expliciter le quasi-inverse.)
